\documentclass[10pt,a4paper]{amsart}
\usepackage{amsmath,amssymb,amsthm,graphicx}
% Retained preamble for the Stacks Project chapter divisors.tex (Divisors).
% The chapter file pinned-original-source.tex carries only \input{preamble} on its first line, so no
% preamble text lives inside the chapter itself.  The Stacks Project uses one single shared preamble for
% every chapter; the snapshot stages it verbatim in sources/stacks-chow/preamble.tex.  The block below is
% copied mechanically, line for line, from that file's lines 14-18 (the Xy-pic setup) and lines 154-217
% (the native theorem declarations together with the mathematical macros).  Nothing was added, removed or
% redefined.  The wrapper class and the amsmath/amssymb/amsthm/graphicx base packages are supplied by the
% assembler, so the class/\externaldocument part of the shared preamble is not repeated here.
% For commutative diagrams we use Xy-pic
\usepackage[all]{xy}

% We use 2cell for 2-commutative diagrams.
\xyoption{2cell}
% Theorem environments.
%
\theoremstyle{plain}
\newtheorem{theorem}[subsection]{Theorem}
\newtheorem{proposition}[subsection]{Proposition}
\newtheorem{lemma}[subsection]{Lemma}

\theoremstyle{definition}
\newtheorem{definition}[subsection]{Definition}
\newtheorem{example}[subsection]{Example}
\newtheorem{exercise}[subsection]{Exercise}
\newtheorem{situation}[subsection]{Situation}

\theoremstyle{remark}
\newtheorem{remark}[subsection]{Remark}
\newtheorem{remarks}[subsection]{Remarks}

\numberwithin{equation}{subsection}

% Macros
%
\def\lim{\mathop{\mathrm{lim}}\nolimits}
\def\colim{\mathop{\mathrm{colim}}\nolimits}
\def\Spec{\mathop{\mathrm{Spec}}}
\def\Hom{\mathop{\mathrm{Hom}}\nolimits}
\def\Ext{\mathop{\mathrm{Ext}}\nolimits}
\def\SheafHom{\mathop{\mathcal{H}\!\mathit{om}}\nolimits}
\def\SheafExt{\mathop{\mathcal{E}\!\mathit{xt}}\nolimits}
\def\Sch{\mathit{Sch}}
\def\Mor{\mathop{\mathrm{Mor}}\nolimits}
\def\Ob{\mathop{\mathrm{Ob}}\nolimits}
\def\Sh{\mathop{\mathit{Sh}}\nolimits}
\def\NL{\mathop{N\!L}\nolimits}
\def\CH{\mathop{\mathrm{CH}}\nolimits}
\def\proetale{{pro\text{-}\acute{e}tale}}
\def\etale{{\acute{e}tale}}
\def\QCoh{\mathit{QCoh}}
\def\Ker{\mathop{\mathrm{Ker}}}
\def\Im{\mathop{\mathrm{Im}}}
\def\Coker{\mathop{\mathrm{Coker}}}
\def\Coim{\mathop{\mathrm{Coim}}}

% Boxtimes
%
\DeclareMathSymbol{\boxtimes}{\mathbin}{AMSa}{"02}

%
% Macros for moduli stacks/spaces
%
\def\QCohstack{\mathcal{QC}\!\mathit{oh}}
\def\Cohstack{\mathcal{C}\!\mathit{oh}}
\def\Spacesstack{\mathcal{S}\!\mathit{paces}}
\def\Quotfunctor{\mathrm{Quot}}
\def\Hilbfunctor{\mathrm{Hilb}}
\def\Curvesstack{\mathcal{C}\!\mathit{urves}}
\def\Polarizedstack{\mathcal{P}\!\mathit{olarized}}
\def\Complexesstack{\mathcal{C}\!\mathit{omplexes}}
% \Pic is the operator that assigns to X its picard group, usage \Pic(X)
% \Picardstack_{X/B} denotes the Picard stack of X over B
% \Picardfunctor_{X/B} denotes the Picard functor of X over B
\def\Pic{\mathop{\mathrm{Pic}}\nolimits}
\def\Picardstack{\mathcal{P}\!\mathit{ic}}
\def\Picardfunctor{\mathrm{Pic}}
\def\Deformationcategory{\mathcal{D}\!\mathit{ef}}

% --- appended 2026-10-12 (Chief): the project own definitions of the reference/slogan/history
% display environments, copied mechanically from the shared Stacks preamble lines 8-12
% (\usepackage{verbatim} plus the three \newenvironment lines, which the chapter preamble keeps
% verbatim from the project preamble). Nothing added, removed or redefined: these environments
% carry no mathematical content and the project itself typesets them as comments. Without these
% three lines a counted statement whose own text uses them fails with Environment slogan undefined.
\usepackage{verbatim}
\newenvironment{reference}{\comment}{\endcomment}
\newenvironment{slogan}{\comment}{\endcomment}
\newenvironment{history}{\comment}{\endcomment}
\usepackage{cleveref}
\crefname{originallemma}{original lemma}{original lemmas}
\Crefname{originallemma}{Original lemma}{Original lemmas}
\makeatletter
\newlabel{algebra-lemma-regular-sequence-in-neighbourhood}{{061L}{source}{}{sourcepointer-1}{}}
\newlabel{algebra-lemma-regular-sequence-in-neighbourhood@cref}{{[originallemma][0][]061L}{[1][0][]source}{}{}{}}
\makeatother
\title{A closed subscheme whose ideal is locally generated by one element, or by a single nonzerodivisor, is locally principal, respectively an effective Cartier divisor}
\author{The Stacks Project Authors}
\date{stacks-project:divisors.tex:local-snapshot-20261012; GNU FDL 1.2 or later}
\begin{document}
\maketitle
\noindent\textit{Editorial note.} All delivered mathematical text is the author's own, copied from the pinned Stacks Project chapter file pinned-original-source.tex; no statement or proof step was rewritten, reordered or removed.\par\medskip
\noindent\textit{Editorial note.} Following the multi-claim convention for this source, source\_url is the source-level URL https://stacks.math.columbia.edu/ and the source label lemma-effective-Cartier-in-points is carried in claim\_key and here rather than in a per-claim tag URL.\par\medskip
\noindent\textit{Editorial note.} The preamble is the chapter's shared Stacks Project preamble, retained as sources/stacks-divisors/retained-preamble.tex: a line-for-line excerpt of the shared Stacks preamble staged with this snapshot. No macro was invented or redefined, and this card adds no preamble text of its own.\par\medskip
\noindent\textit{Editorial note.} Statement span 2938-2952 and proof span 2954-2972 come from the halfA pool card operations/pool-cards/stacks-divisors-halfA.json and were re-read in the pinned chapter file before this card was written.\par\medskip
\noindent\textit{Editorial note.} The proof introduces the finitely many generators of the ideal on the chart and the denominators to be inverted; the author's notation is preserved.\par\medskip
\noindent\textit{Editorial note.} 1 cross-references occur inside the counted spans; they are listed in external\_references as navigation pointers only.\par\medskip
\noindent\textit{Allowed standard prerequisites (editorial).}\begin{itemize}
\item The localization of a finitely generated ideal at a prime is generated by the localizations of its generators, and a principal generator at the prime can be written with denominators outside the prime. Used at proof lines 2954-2967. Produces a neighbourhood on which the ideal becomes principal, giving the first part.
\item An element whose image in a localization at a prime is a nonzerodivisor is a nonzerodivisor on a neighbourhood of the prime, cited from the algebra chapter. Cited at proof lines 2969-2971. Upgrades the pointwise nonzerodivisor condition to a neighbourhood condition, giving the second part.
\end{itemize}\medskip
\ifdefined\hypertarget\hypertarget{sourcepointer-1}{}\fi
\noindent\textit{Source reference.} Original lemma 061L: focused-seeds/AGP/sources/stacks-snapshot-20261011/tags.txt, source lines 7765--7765. Same-chapter reference used in the retained author text at source line(s) 2970, 4250, 4869, 5359, 5534, 5635. It is supplied only as a navigation pointer so that the retained text compiles; the counted item delivers the author's own statement and the author's complete proof, and no part of the referenced item is copied in or substituted for the counted argument.\par\medskip
\begin{lemma}
\label{lemma-effective-Cartier-in-points}
Let $X$ be a locally Noetherian scheme. Let $D \subset X$ be a closed subscheme
corresponding to the quasi-coherent ideal sheaf
$\mathcal{I} \subset \mathcal{O}_X$.
\begin{enumerate}
\item If for every $x \in D$ the ideal
$\mathcal{I}_x \subset \mathcal{O}_{X, x}$
can be generated by one element, then $D$ is locally principal.
\item If for every $x \in D$ the ideal
$\mathcal{I}_x \subset \mathcal{O}_{X, x}$
can be generated by a single nonzerodivisor, then $D$ is an
effective Cartier divisor.
\end{enumerate}
\end{lemma}

\begin{proof}
Let $\Spec(A)$ be an affine neighbourhood of a point $x \in D$.
Let $\mathfrak p \subset A$ be the prime corresponding to $x$.
Let $I \subset A$ be the ideal defining the trace of $D$ on
$\Spec(A)$. Since $A$ is Noetherian (as $X$ is locally Noetherian)
the ideal $I$ is generated by finitely many elements, say
$I = (f_1, \ldots, f_r)$. Under the assumption of (1) we have
$I_\mathfrak p = (f)$ for some $f \in A_\mathfrak p$.
Then $f_i = g_i f$ for some $g_i \in A_\mathfrak p$.
Write $g_i = a_i/h_i$ and $f = f'/h$ for some
$a_i, h_i, f', h \in A$, $h_i, h \not \in \mathfrak p$.
Then $I_{h_1 \ldots h_r h} \subset A_{h_1 \ldots h_r h}$ is
principal, because it is generated by $f'$. This proves (1).
For (2) we may assume $I = (f)$. The assumption implies
that the image of $f$ in $A_\mathfrak p$ is a nonzerodivisor.
Then $f$ is a nonzerodivisor on a neighbourhood of $x$ by
Algebra, Lemma \ref{algebra-lemma-regular-sequence-in-neighbourhood}.
This proves (2).
\end{proof}

\end{document}
